\documentclass[11pt]{article}

\usepackage[margin=2.5cm]{geometry}
\usepackage{amsmath,amssymb}
\usepackage{hyperref}

\newcommand{\pd}[2]{\frac{\partial #1}{\partial #2}}

\title{Derivations}
\author{Matthew Archer}
\date{}

\begin{document}
\maketitle
\tableofcontents

% =====================================================================
\section{Single-example MLP (\texttt{torch\_linear.ipynb})}
% =====================================================================

\subsection{Explanation of the Jacobian check}

We want to be able to build the jacobian numerically and analytically in order
to check the derivatives.

Analytically, we can use the vjp and use a unit vector to retrieve the jacobian
row by row.

Numerically, we can compute finite differences to retrieve the jacobian column
by column.

Example:

Definitions: $W$ is (rows, cols), $x$ is a column vector with (cols, 1)

For the $\frac{dy}{dW}$ (i.e.\ local) Matrix vector multiplication: $y = W x$

The jacobain $dy / dW$ is a rank 3 tensor:

Numerically we loop over $W$ (row by column). For each index $W_{ij}$ we get a
column vector of shape (cols, 1). When stacked this is (8, 2, 1)

Analytically, we recover a `row' for each output $y_i$. In this case, a row
turns out to contain a matrix, this has shape $[2, 2, 4]$.

\medskip
Alternatively we can see this as the following:

\texttt{J\_analytic = sens\_W\_analytic.view(rows, rows * cols)}

i.e.\ this flattens the (2, 4, 2) into (2, 8)

Our numeric is already (8, 2, 1)

Which we could squeeze and take the transpose to get (2, 8)

\subsection{Derivation of vjp for mm, add..}

\subsubsection*{Forward}

Layer 1
\begin{align*}
s_i &= W_{ij} x_j \\
u_i &= s_i + b_i
\end{align*}

Activation
\[
z_i = \sigma(u_i)
\]

Layer 2
\begin{align*}
t_i &= B_{ij} z_j \\
\hat{y}_i &= t_i + c_i
\end{align*}

Loss
\[
l = 0.5\,(\hat{y}_i - y_i)^2
\]

\subsubsection*{Backward}
Sequence of steps:

Keeping the outer (upstream) gradient on the left. (This doesn't really matter
when using index notation)

Error (on a single example basis)
\[
\frac{dl}{d\hat{y_i}} = \hat{y_i} - y_i
\]

2nd Add, $\hat{y}_i = t_i + c_i$
\[
\pd{l}{t_i} = \pd{l}{\hat{y_k}} \pd{\hat{y_k}}{t_i}, \quad \pd{l}{c_i} = \pd{l}{\hat{y_k}} \pd{\hat{y_k}}{c_i}
\]

2nd Matmul, $t_i = B_{ij} z_j$
\[
\pd{l}{B_{ij}} = \pd{l}{t_k} \pd{t_k}{B_{ij}}, \quad \pd{l}{z_i} = \pd{l}{t_k} \pd{t_k}{z_i}
\]

Propagate through activation function, $z_i = \sigma(u_i)$
\[
\pd{l}{u_i} = \pd{l}{z_k}\pd{z_k}{u_i}
\]

1st Add, $u_i = s_i + b_i$
\[
\pd{l}{s_i} = \pd{l}{u_k} \pd{u_k}{s_i} , \quad \pd{l}{b_i} = \pd{l}{u_k} \pd{u_k}{b_i}
\]

1st Matmul, $s_i = W_{ij} x_j$
\[
\pd{l}{W_{ij}} = \pd{l}{s_k}\pd{s_k}{W_{ij}}, \quad \pd{l}{x} \text{ is unused}
\]

\subsection{Actual derivations assuming scalar input}
First add is trivial so we'll ignore. Error is actually $e_i = \hat{y_i} - y_i$

Sum is $\Sigma_i^n$, where $n$ is the batch size. $l = 0.5e_i e_i$. i.e.
\[
\frac{dl}{d\hat{y_i}} = \hat{y_i} - y_i = e_i
\]

2nd Add, $\hat{y}_i = t_i + c_i$
\[
\pd{l}{t_i} = \pd{l}{\hat{y_k}} \pd{\hat{y_k}}{t_i}, \quad \pd{l}{c_i} = \pd{l}{\hat{y_k}} \pd{\hat{y_k}}{c_i}
\]
\[
\pd{l}{t_i} = e_k \delta_{ki}, \quad \pd{l}{c_i} = e_k \delta_{ki}
\]

2nd Matmul, $t_i = B_{ij} z_j$
\[
\pd{l}{B_{ij}} = \pd{l}{t_k} \pd{t_k}{B_{ij}}, \quad \pd{l}{z_i} = \pd{l}{t_k} \pd{t_k}{z_i}
\]
\[
\pd{l}{B_{ij}} = e_k \pd{B_{kl} z_l}{B_{ij}} = e_k \delta_{ki} \delta_{lj} z_l = e_i z_j, \quad \pd{l}{z_i} = e_k \pd{B_{kl} z_l}{z_i} = e_k B_{kl} \delta_{li} = e_k B_{ki}
\]

Propagate through activation function, $z_i = \sigma(u_i)$
\[
\pd{l}{u_i} = \pd{l}{z_k}\pd{z_k}{u_i}
\]
\[
\pd{l}{u_i} = e_m B_{mk} \pd{\sigma(u_k)}{u_i} = e_m B_{mk} \sigma_k(1-\sigma_k)\delta_{ki} = e_m B_{mi} \sigma_i(1-\sigma_i)
\]

1st Add, $u_i = s_i + b_i$
\[
\pd{l}{s_i} = \pd{l}{u_k} \pd{u_k}{s_i} , \quad \pd{l}{b_i} = \pd{l}{u_k} \pd{u_k}{b_i}
\]

Addition just streams upstream derivatives (adjoint)
\[
\pd{l}{s_i} = \pd{l}{u_k} \delta_{ki} = \pd{l}{u_i}, \quad \pd{l}{b_i} = \pd{l}{u_i}
\]

1st Matmul, $s_i = W_{ij} x_j$
\[
\pd{l}{W_{ij}} = \pd{l}{s_k}\pd{s_k}{W_{ij}}, \quad \pd{l}{x} \text{ is unused}
\]
\[
\pd{l}{W_{ij}} = \pd{l}{u_k}\pd{W_{kl} x_l}{W_{ij}} = \pd{l}{u_k}\delta_{ki} \delta_{lj} x_l = \pd{l}{u_i} x_j
\]

TODO: I don't think this works for batches yet as it assumes that $x_j$ is just
a vector -- need to write it for batch next.

% =====================================================================
\section{Batch implementation (\texttt{torch\_linear\_batch\_derivation.ipynb})}
% =====================================================================

This is a reimplementation of \texttt{torch\_linear.ipynb} computing SGD on
batches. The mathematics isn't really that much more difficult.

\subsection{Derivation of vjp for mm, add..}

\subsubsection*{Forward}

Layer 1

We typically treat $W$ as (out\_dim, in\_dim), which is fine when $x$ is a
vector. However, with batched inputs, we want each row to correspond to a
different item in the batch: $X_{ik}$ -- each row is an example -- $i$
increments over the batch.

Previously we had.
\begin{equation*}
s_{i} = W_{ij} x_j
\end{equation*}

Now batch aware ($i$ is the batch index):
\begin{align*}
  s_{ij} &=  W_{jk} X_{ik} \\
  s_{ij} &=  W^T_{kj} X_{ik} \\
  s_{ij} &=  X_{ik} W^T_{kj}
\end{align*}

In non-index notation, this is:
\begin{equation*}
  \mathbf{s}  := \mathbf{X} @ \mathbf{W}^T
\end{equation*}
Broadcasting:
\begin{equation*}
u_{ij} = s_{ij} + b_{j}
\end{equation*}
Activation
\[
z_{ij} = \sigma(u_{ij})
\]

Layer 2

$B$ is actually a vector for univariate regression. But we write it in general
multivariate here.
\begin{align*}
t_{ij} &= z_{ik} B^T_{kj} \\
\hat{y}_{ij} &= t_{ij} + c_j
\end{align*}

Loss (where $i$ is the example and $j$ is the number of outputs). Remember that
$y$ is the actual value and $\hat{y}$ is the prediction. Below it is written in
general form for multi variable output. I.e.\ regression in multiple dimensions
(predicting 2 dimensions for instance.)
\[
l = 0.5\,(\hat{y}_{ij} - y_{ij})(\hat{y}_{ij} - y_{ij})
\]

\subsubsection*{Backward}
Sequence of steps:

Keeping the outer (upstream) gradient on the left. (This doesn't really matter
when using index notation)
\[
\frac{dl}{d\hat{y}_{kl}} = \frac{1}{2NM} \pd{}{\hat{y}_{kl}}(\hat{y}_{ij} - y_{ij})^2 = \frac{1}{NM} \delta_{ik}\delta_{jl} (\hat{y}_{ij} - y_{ij}) = \frac{1}{NM}(\hat{y}_{kl} - y_{kl}) = e_{kl}
\]

2nd Add, $\hat{y}_{ij} = t_{ij} + c_j$
\[
\pd{l}{t_{ij}} = \pd{l}{\hat{y}_{kl}} \pd{\hat{y}_{kl}}{t_{ij}}, \quad \pd{l}{c_j} = \pd{l}{\hat{y}_{kl}} \pd{\hat{y}_{kl}}{c_j}
\]
\[
\pd{l}{\hat{y}_{kl}} = e_{kl}
\]
\[
\pd{l}{t_{ij}} = e_{ij}, \quad \pd{l}{c_j} = \frac{1}{NM} (\hat{y}_{kj} - y_{kj}) = \sum_k^N  e_{kj}
\]

2nd Matmul, $t_{ij} = z_{ik} B^{T}_{kj}$
\[
\pd{l}{B_{ij}} = \pd{l}{t_{kl}} \pd{t_{kl}}{B_{ij}}, \quad \pd{l}{z_{ij}} = \pd{l}{t_{kl}} \pd{t_{kl}}{z_{ij}}
\]
\[
\pd{l}{t_{kl}} = e_{kl}, \quad \pd{t_{kl}}{B_{ij}} = z_{km}\pd{B_{ml}^T}{B_{ij}} = z_{km}\delta_{il}\delta_{jm}
\]
\[
\pd{l}{B_{ij}} = e_{ki} z_{kj}, \quad \pd{l}{z_{ij}} = e_{il}B_{lj}
\]

Propagate through activation function, $z_{ij} = \sigma(u_{ij})$
\[
\pd{l}{u_{ij}} = \pd{l}{z_{kl}}\pd{z_{kl}}{u_{ij}}
\]
\[
\pd{l}{u_{ij}} = e_{km}{B_{ml}}\sigma_{kl}(1-\sigma_{kl})\delta_{ik}\delta_{jl} = e_{im}{B_{mj}}\sigma_{ij}(1-\sigma_{ij})
\]

1st Add, $u_{ij} = s_{ij} + b_j$
\[
\pd{l}{s_{ij}} = \pd{l}{u_{kl}} \pd{u_{kl}}{s_{ij}} , \quad \pd{l}{b_j} = \pd{l}{u_{kl}} \pd{u_{kl}}{b_j}
\]
\[
\pd{l}{s_{ij}} = \pd{l}{u_{kl}} \delta_{ik}\delta_{jl}, \quad \pd{l}{b_j} = \pd{l}{u_{kl}} \delta_{jl}
\]

1st Matmul, $s_{ij} = X_{ik} W_{kj}^T$
\[
\pd{l}{W_{ij}} = \pd{l}{s_{kl}}\pd{s_{kl}}{W_{ij}}, \quad \pd{l}{X_{ij}} \text{ is unused}
\]
\[
\pd{l}{W_{ij}} = \pd{l}{s_{kl}} X_{km} \pd{W_{ml}^T}{W_{ij}} = \pd{l}{s_{kl}} X_{km} \delta_{il} \delta_{jm} = \pd{l}{s_{ki}} X_{kj}
\]

\subsection{Test problem: Rotations}

Given $\mathbf{x} \in \mathbb{R}^2$ and $\theta \in (0, 2 \pi)$ we must learn
the 2D rotation matrix:
\[
\mathbf{y} = \mathbf{R} \mathbf{x} = \begin{pmatrix} cos\theta & -sin\theta \\ sin\theta & cos\theta \end{pmatrix} \begin{pmatrix} x_1 \\ x_2 \end{pmatrix}
\]

% =====================================================================
\section{Batch norm derivation (\texttt{torch\_batch\_norm\_derivation.ipynb})}
% =====================================================================

\[
y = \frac{x - \mathrm{E}[x]}{\sqrt{\mathrm{Var}[x] + \epsilon}} * \gamma + \beta
\]

Best pen and paper source is:

Kevin Zakka's blog post: ``Deriving the Gradient for the Backward Pass of Batch
Normalization''.

\begin{enumerate}
  \setcounter{enumi}{3}
  \item The Original Source: Batch Normalization: Accelerating Deep Network
        Training by Reducing Internal Covariate Shift
\end{enumerate}
\begin{enumerate}
  \setcounter{enumi}{2}
  \item The Math Derivation: Deriving the Gradient for the Backward Pass of
        Batch Normalization. By: Frederik Kratzert (Blog Post / Markdown Guide)
\end{enumerate}

\url{https://www.youtube.com/watch?v=Y23QgQGAGJQ}

Matrix calculus you need for DL. \url{https://arxiv.org/pdf/1802.01528}

\begin{enumerate}
  \item It explicitly focuses on the Vector-Jacobian Product (VJP). Instead of
        forcing you to memorize giant Jacobian matrices, Howard and Parr
        introduce the math from a ``row vector times Jacobian'' perspective.
        They show you how to trace a vector of downstream gradients through
        operations without ever materializing a giant, sparse, memory-wasting
        matrix in your code.
\end{enumerate}

Strategy:
\begin{itemize}
  \item decompose all the above into individual vjp
\end{itemize}

\subsection{Now write in terms of ops}

\begin{enumerate}
  \setcounter{enumi}{-1}
  \item Summation / scaling gate (mean) -- $1/n$ sum over batch
  \item Subtraction
  \item Power
  \item Summation / scaling ()
  \item Denominator
  \item inverse
  \item mm
\end{enumerate}

\subsubsection*{Example}

$X$ is (3, 5) i.e.\ 3 examples and 5 features

mean should be size 5 (average over each feature)

\subsection{Now do vjps}

\subsubsection*{Vjp of broadcast and reduction}
To make this easier I have explicitly used broadcast and reduction in the
forward chain. I also broke the sum scale into two separate ops. Scaling by a
constant is trivial.
\begin{itemize}
  \item there are some downsides to this that I will discuss later
  \item benefits
\end{itemize}
However we have to put in control flow checks for different dimensions of a
tensor.

some notes:
\begin{itemize}
  \item add\_scalar is trivial
\end{itemize}

We now need to implement these vjps in reverse.

\subsubsection*{vjp of \texttt{broadcast\_to()}}
The vjp of the broadcast

Going from $b_j$ to $B_{ij}$

$B_{ij} = \mathrm{fn}(b_j)$ therefore we differentiated $L$ w.r.t.\ $B_{ij}$

\paragraph{Broadcast.} Forward:
\[
B_{ij} = b_j
\]
vjp of broadcast:
\[
\pd{L}{b_k} = \pd{L}{B_{ij}} \pd{B_{ij}}{b_k} = \pd{L}{B_{ij}}\delta_{kj} = \sum_i \pd{L}{B_{ik}}
\]

\paragraph{Reduction.} Forward:
\[
b_j = \sum_i B_{ij}
\]
vjp of reduction:
\[
\pd{L}{B_{ij}} = \pd{L}{b_k} \pd{b_{k}}{B_{ij}} = \pd{L}{b_k}\sum_l \pd{B_{lk}}{B_{ij}} = \pd{L}{b_k}\sum_l \delta_{il} \delta_{jk} = \pd{L}{b_j}
\]
i.e.\ the r.h.s sum is non-zero only when $l = i$.

So we see that the vjp of a broadcast is a reduction, and the vjp of a sum is a
broadcast.

\subsubsection*{vjp of \texttt{hadamard()}}
\begin{align*}
p &= a \odot b\\
\pd{L}{a_{ij}} &= \pd{L}{p_{kl}} \pd{p_{kl}}{a_{ij}} \\
\pd{L}{b_{ij}} &= \pd{L}{p_{kl}} \pd{p_{kl}}{b_{ij}} \\
p_{ij} &= a_{ij} b_{ij}\\
\pd{p_{ij}}{a_{kl}} = \pd{(a_{ij} b_{ij})}{a_{kl}} = b_{ij} \pd{a_{ij}}{a_{kl}} &= b_{ij}\delta_{ki}\delta_{lj} \\
 \pd{p_{kl}}{a_{ij}}  &= b_{kl}\delta_{ik}\delta_{jl} \\
\pd{L}{a_{ij}} &= \pd{L}{p_{kl}} b_{kl} \delta_{ik}\delta_{jl} = \pd{L}{p_{ij}} b_{ij}
\end{align*}

ops: $dL/dB$
\[
B = B(\gamma, \beta)
\]
\[
s = \sum_b X_{bj}
\]

vjp
\[
\frac{\partial s }{\partial X_{lm}} = \sum_b \frac{\partial X}{\partial}
\]

\subsection{Other ops?}
\begin{itemize}
  \item Broadcasting might be an option to be absolutely explicit -- we have to
        do the opposite for vjp.
\end{itemize}

\end{document}
